
\documentclass{article}
\usepackage{colortbl}
\usepackage{makecell}
\usepackage{multirow}
\usepackage{supertabular}

\begin{document}

\newcounter{utterance}

\centering \large Interaction Transcript for game `cladder', experiment `full\_v1.5\_default', episode 4422 with refineonly{-}llama{-}3.1{-}8b{-}instruct.
\vspace{24pt}

{ \footnotesize  \setcounter{utterance}{1}
\setlength{\tabcolsep}{0pt}
\begin{supertabular}{c@{$\;$}|p{.15\linewidth}@{}p{.15\linewidth}p{.15\linewidth}p{.15\linewidth}p{.15\linewidth}p{.15\linewidth}}
    \# & \multicolumn{2}{c}{Player} && \multicolumn{2}{c}{Game Master} \\
    \hline

    \theutterance \stepcounter{utterance}  
    & & & \multicolumn{4}{p{0.6\linewidth}}{
        \cellcolor[rgb]{0.9,0.9,0.9}{
            \makecell[{{p{\linewidth}}}]{
                \texttt{\tiny{[P1$\langle$GM]}}
                \texttt{You are an expert in causal inference. The following question is not a typical commonsense query, but rather a meticulously designed question created by a professor specializing in causal inference, intended to assess the students' mastery of the course content.} \\
\texttt{Imagine a self{-}contained, hypothetical world with only the following conditions, and without any unmentioned factors or causal relationships: Gender has a direct effect on smoking and peanut allergy. Smoking has a direct effect on tar deposit. Tar deposit has a direct effect on peanut allergy. Gender is unobserved. For nonsmokers, the probability of high tar deposit is 18\%. For smokers, the probability of high tar deposit is 64\%. For nonsmokers and with no tar deposit, the probability of being allergic to peanuts is 46\%. For nonsmokers and with high tar deposit, the probability of being allergic to peanuts is 77\%. For smokers and with no tar deposit, the probability of being allergic to peanuts is 17\%. For smokers and with high tar deposit, the probability of being allergic to peanuts is 43\%. The overall probability of smoking is 49\%. Will smoking decrease the chance of being allergic to peanuts?} \\
\texttt{Start your answer with 'Yes' or 'No', followed by additional reasoning or evidence to support your explanation.} \\
            }
        }
    }
     \\ \\

    \theutterance \stepcounter{utterance}  
    & \multicolumn{4}{p{0.6\linewidth}}{
        \cellcolor[rgb]{0.561,0.737,0.561}{
            \makecell[{{p{\linewidth}}}]{
                \texttt{\tiny{[P1$\rangle$GM]}}
                \texttt{No, smoking will not decrease the chance of being allergic to peanuts.} \\
\\ 
\texttt{To answer this question, we need to use the concept of counterfactuals and potential outcomes. We are given the following probabilities:} \\
\\ 
\texttt{{-} P(high tar deposit | smoking) = 0.64} \\
\texttt{{-} P(high tar deposit | nonsmoking) = 0.18} \\
\texttt{{-} P(allergic to peanuts | nonsmoking, high tar deposit) = 0.77} \\
\texttt{{-} P(allergic to peanuts | nonsmoking, no tar deposit) = 0.46} \\
\texttt{{-} P(allergic to peanuts | smoking, high tar deposit) = 0.43} \\
\texttt{{-} P(allergic to peanuts | smoking, no tar deposit) = 0.17} \\
\texttt{{-} P(smoking) = 0.49} \\
\\ 
\texttt{We want to compare the probability of being allergic to peanuts among smokers and nonsmokers. To do this, we need to calculate the probability of being allergic to peanuts among nonsmokers and smokers, respectively.} \\
\\ 
\texttt{First, let's calculate the probability of being allergic to peanuts among nonsmokers. We can use the law of total probability:} \\
\\ 
\texttt{P(allergic to peanuts | nonsmoking) = P(allergic to peanuts | nonsmoking, high tar deposit) * P(high tar deposit | nonsmoking) + P(allergic to peanuts | nonsmoking, no tar deposit) * P(no tar deposit | nonsmoking)} \\
\\ 
\texttt{We know that P(high tar deposit | nonsmoking) = 0.18 and P(allergic to peanuts | nonsmoking, high tar deposit) = 0.77. To find P(no tar deposit | nonsmoking), we can use the fact that P(high tar deposit | nonsmoking) + P(no tar deposit | nonsmoking) = 1. Therefore, P(no tar deposit | nonsmoking) = 1 {-} 0.18 = 0.82.} \\
\\ 
\texttt{Now we can plug in the values:} \\
\\ 
\texttt{P(allergic to peanuts | nonsmoking) = 0.77 * 0.18 + 0.46 * 0.82} \\
\texttt{= 0.1391 + 0.3772} \\
\texttt{= 0.5163} \\
\\ 
\texttt{Next, let's calculate the probability of being allergic to peanuts among smokers. We can use the law of total probability:} \\
\\ 
\texttt{P(allergic to peanuts | smoking) = P(allergic to peanuts | smoking, high tar deposit) * P(high tar deposit | smoking) + P(allergic to peanuts | smoking, no tar deposit) * P(no tar deposit | smoking)} \\
\\ 
\texttt{We know that P(high tar deposit | smoking) = 0.64 and P(allergic to peanuts | smoking, high tar deposit) = 0.43. To find P(no tar deposit | smoking), we can use the fact that P(high tar deposit | smoking) + P(no tar deposit | smoking) = 1. Therefore, P(no tar deposit} \\
            }
        }
    }
    & & \\ \\

    \theutterance \stepcounter{utterance}  
    & & & \multicolumn{2}{p{0.3\linewidth}}{
        \cellcolor[rgb]{0.9,0.9,0.9}{
            \makecell[{{p{\linewidth}}}]{
                \texttt{\tiny{[GM$|$GM]}}
                \texttt{No} \\
            }
        }
    }
    & & \\ \\

    \theutterance \stepcounter{utterance}  
    & & & \multicolumn{2}{p{0.3\linewidth}}{
        \cellcolor[rgb]{0.9,0.9,0.9}{
            \makecell[{{p{\linewidth}}}]{
                \texttt{\tiny{[GM$|$GM]}}
                \texttt{no} \\
            }
        }
    }
    & & \\ \\

    \theutterance \stepcounter{utterance}  
    & & & \multicolumn{2}{p{0.3\linewidth}}{
        \cellcolor[rgb]{0.9,0.9,0.9}{
            \makecell[{{p{\linewidth}}}]{
                \texttt{\tiny{[GM$|$GM]}}
                \texttt{game\_result = WIN} \\
            }
        }
    }
    & & \\ \\

\end{supertabular}
}

\end{document}
